\documentclass[12pt]{article}
\usepackage[utf8]{inputenc}
\usepackage[T1]{fontenc}
\usepackage{hyperref}
\title{Uisai Language Resources for K-12 Curriculum Development}
\author{Compiled by Claude Sonnet 5 (claude-sonnet-5) for Andr\'as Kornai}
\date{}
\begin{document}
\maketitle

\section{Introduction}
Uisai (ISO 639-3 uis) is a South Bougainville language of the Buin area of Bougainville, Papua New Guinea. The endangered-languages database cites about 2,500 speakers (1991 SIL data) and rates it threatened. Very little material was found. Note that the Joshua Project profile places Uisai on New Britain, which conflicts with other sources and is treated here as unreliable.

\section{Primary Educational Resources}
No primers, textbooks or school programmes were found in the search results.

\section{Linguistic and Reference Materials}
Overview: \url{https://en.wikipedia.org/wiki/Uisai_language}. Endangered Languages Project entry: \url{https://www.endangeredlanguages.com/lang/5845}. Paradisec holds 1960 recordings and a vocabulary list by Arthur Capell for an inland Buin dialect: \url{https://paradisec.org.au/fieldnotes/NSBUIS.htm}.

\section{Religious and Community Resources}
None found.

\section{Digital and Audio Resources}
Paradisec archive item: \url{https://catalog.paradisec.org.au/collections/AC1/items/002}.

\section{Curriculum Development Context}
Although the language is described as used in education, no curriculum documents were found. Development would need to start from field contact and archival materials such as the Paradisec holdings.
\end{document}